% Section 4.3: LoRA 微调方案设计

LoRA 分支在 Base 之上引入训练侧参数化，覆盖 4.1 节矩阵的 H4--H10 七组。三种适配器均在同一基座、同一训练脚本骨架下产出，差异仅在训练样本的监督目标与是否在训练 prompt 中注入 few-shot 示例。本节 4.3.1 先界定三种适配器的训练样本与 loss 口径，4.3.2 给出与之匹配的一组超参数，4.3.3 说明从训练落盘到推理侧挂载的完整链路。

\subsection{训练样本与目标}
% Subsection 4.3.1: 训练样本与目标

三种 LoRA 适配器（\texttt{LoRA\_code}、\texttt{LoRA\_cot\_zs}、\texttt{LoRA\_cot\_fs(k)}）共用 \texttt{KodCode\_train.jsonl}，但在 assistant 回复的构造与 loss 计算口径上有所区分，以匹配各自对应的推理场景。

\subsubsection{三类监督目标}
同一训练集经不同脚本处理后形成三种监督目标（见表~\ref{tab:ch4-lora-target}），支撑 H4/H7/H9 三条训练路径；具体脚本与指针文件对应 4.1 节表~\ref{tab:ch4-matrix}。
\begin{table}[!htb]
  \centering
  \caption{三种 LoRA 适配器的监督目标与训练 prompt 配置}
  \label{tab:ch4-lora-target}
  \footnotesize
  \setlength{\tabcolsep}{5pt}
  \renewcommand{\arraystretch}{1.2}
  \begin{tabular}{c l c p{5cm}}
    \toprule
    适配器 & assistant 回复内容 & few-shot & 产出脚本 \\
    \midrule
    LoRA\_code  &  solution（纯代码块）  &  无  &  train/train.py \\
    LoRA\_cot\_zs  &  thought\_step + solution  &  0 例  &  train\_cot.py --num-few-shots 0 \\
    LoRA\_cot\_fs($k$)  &  thought\_step + solution  &  $k$ 例  &  train\_cot.py --num-few-shots k \\
    \bottomrule
  \end{tabular}
\end{table}


\subsubsection{Chat 样本结构}
\begin{itemize}
    \item \textbf{消息条数}：\texttt{LoRA\_code} / \texttt{LoRA\_cot\_zs} 为 system + user + assistant 共 3 条；\texttt{LoRA\_cot\_fs($k$)} 为 $2{+}2k{+}2$ 条多轮消息（$k$ 组示范 + 当前题）；
    \item \textbf{角色分工}：system 固化角色；user 承载题面、函数签名与指令；assistant 承载监督目标；
    \item \textbf{字面模板}：各档在 \texttt{apply\_chat\_template} 之前的完整骨架见 4.1 节（1a）；
    \item \textbf{相对 Direct 的 CoT 差异}：user 末尾追加 \texttt{Think step by step, then output the Python code.}；assistant 为「推理段 + 空行 + fenced Python」，与 4.4 节评测侧对 few-shot 示例的回放形态一致。
\end{itemize}

\subsubsection{Loss 掩码策略}
为避免在 system/user 段上消耗模型容量或干扰指令跟随，loss 仅在最后一条 assistant 回复上反向传播：通过对完整序列与去除目标回复的前缀进行对齐定位 assistant 起点，将其之前的位置全部置为忽略标记。该实现沿用基座自带的 chat template 来解释特殊 token，避免手工拼接字符串带来的 tokenizer 不一致。序列长度上限按 LoRA\_code 与 LoRA\_cot\_* 两类适配器分别设置，以兼顾 CoT 监督目标段的完整保留与单卡显存预算。

\subsubsection{Few-shot 训练样本的构造}
\texttt{LoRA\_cot\_fs($k$)} 独有环节：对每条目标样本 $q$，用 3.3.2 节定义的强标签集合派生 \texttt{strong\_tags}，在训练池中按下式评分：
\[
\text{score}(q, q') = 2\cdot |\mathrm{strong\_tags}(q) \cap \mathrm{strong\_tags}(q')| + |\mathrm{tags}(q) \cap \mathrm{tags}(q')|
\]
\begin{itemize}
    \item \textbf{候选池}：取得分最高的前 $\max(k, 4k)$ 条；
    \item \textbf{抽样}：\texttt{random.sample} 从中抽 $k$ 条，以 User$\rightarrow$Assistant 轮次插在目标题之前，形成 \texttt{(题目, 推理, 代码)} 示范序列；
    \item \textbf{与评测对齐}：与 4.4.2 节评测侧检索共享同一评分与强标签集合；
    \item \textbf{随机性差异}：训练为全局 RNG；评测侧 seed 由题号与采样编号确定性派生，保证逐条可复现；
    \item \textbf{目的}：train/test 在上下文统计上对齐，同时评测结果可被第三方复核。
\end{itemize}

\subsection{超参数与训练策略}
% Subsection 4.3.2: 超参数与训练策略

三种适配器共用同一套超参数，仅在样本构造层面不同（4.3.1 节）。LoRA 结构上取秩 $r=32$、缩放系数 $\alpha=64$（有效学习率 $\alpha/r=2$），并将适配层注入 ChatGLM 块中的 attention（融合 QKV 与输出）与 MLP 两侧线性层共四类位置，使 CoT 推理段与代码段在参数侧均可获得适配空间——这一点在 \texttt{LoRA\_cot\_*} 的监督目标下尤为关键。优化器侧采用 AdamW 配 cosine 调度（含约 3\% 步的线性 warmup）、bf16 数值精度与 gradient\_checkpointing 以匹配 9B 基座 + 32GB 单卡的显存预算；训练 3 个 epoch、等效 batch size 为 32，与训练集规模（约 9000 条）匹配，足以让 LoRA 收敛而不在 CoT 监督目标上过拟合。具体超参取值在仓库 \texttt{SETUP.md} 中以表格形式给出。

\subsubsection{设计依据与一致性约束}
上述取值遵循以下三条原则，使得 10 组实验之间仅因训练样本构造差异而产生可解释的指标变化：
\begin{itemize}
    \item \textbf{跨适配器一致}：三种 LoRA 共享完全相同的结构与优化器超参，仅在样本构造与示例条数 $k$ 上有差异；
    \item \textbf{与推理侧一致}：\texttt{LoRA\_code} 路径的 chat 模板与 4.2.1 节 Direct 推理对齐；\texttt{LoRA\_cot\_*} 路径的用户指令与 assistant 形态与 4.4 节 CoT 推理（含 few-shot 回放）对齐，避免特殊 token 或围栏格式在 train/test 间错位；
    \item \textbf{与数据规模匹配}：3 epoch 与等效 batch size 在 KodCode 训练集上约合数百次 optimizer step，足以让 LoRA 收敛而不至于过拟合。
\end{itemize}


\subsection{适配器推理流程}
% Subsection 4.3.3: 适配器推理流程

训练脚本将 LoRA 适配器落盘后，由评测端 vLLM 在加载基座的同时挂载相应适配器，覆盖 H4、H7、H9 三组 Direct 推理实验；不同组别之间仅适配器选择不同，挂载方式与具体实现无关。

为保证训练侧主效应可被干净归因，单阶段生成的 Prompt 模板、采样配置、代码抽取规则、子进程执行与超时机制、pass@k 估计公式均与 4.2 节 Base 评测完全一致。结果统一写入 \texttt{evaluation/results/} 下，文件名按 4.1 节表~\ref{tab:ch4-matrix} 末列约定，为第五章差分分析提供文件级可追溯性。
